% Einheit 2 von 3 – Kreisfläche (bank/kreis/e2.jsonl, zusammenbau v0.1)
%% TODO Zweigzeile Teil 1: was hier gelernt wird (Satz in Schülersprache aus dem Einheitstitel) – Einheit 2
\einheitenkopf[e2]{Einheit 2 von 3 · Kreisfläche}
\zweigzeile{ab Kl. 7, je nach Buch bis Kl. 8 · P10 oft}
\setcounter{aufgabe}{14}

%% TODO Titel: Ich-kann-Satz fehlt in der Bank (Platzhalter: Kettenname) – Nr. 15
%% TODO Anweisung: ein Satz über der ersten Teilaufgabe fehlt in der Bank – Nr. 15
\begin{aufgabe}{Fläche}
%% TODO \rechenplatz{n}: Schrittzahl des Grundfalls fehlt in der Bank (nur bei fester Schreibform, 2.3 b)
\begin{teile}
\teil Wie viel Stoff braucht man für eine runde Tischdecke? Welche Formel brauchst du? Kreuze an.\\ \kreuz{$u = \pi \cdot d$}\\ \kreuz{$A = \pi \cdot r^2$}
\teil Kreis mit $r = 13$ cm – wie groß ist die Fläche? Rechne mit der $\pi$-Taste und runde auf eine Stelle. A ≈ \leerfeld[cm²]
\teil Kreis mit $d = 16$ cm – wie groß ist die Fläche? Rechne mit der $\pi$-Taste und runde auf eine Stelle. A ≈ \leerfeld[cm²]
\teil Kreis mit $r = 2{,}6$ m – wie groß ist die Fläche? Rechne mit der $\pi$-Taste und runde auf zwei Stellen. A ≈ \leerfeld[m²]
\teil Eine Bühne hat die Form eines Halbkreises mit $r = 6$ m. Wie groß ist ihre Fläche? Rechne mit der $\pi$-Taste und runde auf eine Stelle. A ≈ \leerfeld[m²]
\teil Der Radius ist gegeben. Ergänze die Tabelle. Rechne mit der $\pi$-Taste und runde auf eine Stelle.

\sachtabelle{cccc}{$r$ & $d$ & $u$ & $A$}{$3{,}5$ cm & \leerzelle & \leerzelle & \leerzelle}
\teil Ein Kreis hat die Fläche $A = 154$ cm². Wie groß ist sein Radius? Rechne mit der $\pi$-Taste und runde auf eine Stelle. r ≈ \leerfeld[cm]
\teil Ein runder Sandkasten hat den Durchmesser $2{,}36$ m. Berechne seinen Flächeninhalt. Runde auf zwei Stellen \hfill (P10 2016 OS). A ≈ \leerfeld
\end{teile}
\end{aufgabe}

%% TODO Titel: Ich-kann-Satz fehlt in der Bank (Platzhalter: Kettenname) – Nr. 16
%% TODO Anweisung: ein Satz über der ersten Teilaufgabe fehlt in der Bank – Nr. 16
\begin{aufgabe}{Term zu Figur}
\begin{teile}
\teil Die Figur ist ein Viertelkreis mit dem Radius $7$ cm. Schreibe einen Term für ihren Flächeninhalt und rechne ihn aus; runde auf eine Stelle.

\begin{kreis}[2] \mittelpunkt{M} \sektor{0}{90}{} \radius{0}{$7$ cm} \end{kreis}
Term: \leerfeld ≈ \leerfeld[cm²]
\end{teile}
\end{aufgabe}

%% TODO Titel: Ich-kann-Satz fehlt in der Bank (Platzhalter: Kettenname) – Nr. 17
%% TODO Anweisung: ein Satz über der ersten Teilaufgabe fehlt in der Bank – Nr. 17
\begin{aufgabe}{Sachaufgabe (Abwurfring, Pizza, Deckel, Grundfläche eines Kegels oder Zylinders)}
\begin{teile}
\teil Der runde Deckel eines Eimers hat den Durchmesser $28$ cm. Wie groß ist seine Fläche? Rechne mit der $\pi$-Taste und runde auf ganze cm². ≈ \leerfeld[cm²]
\end{teile}
\end{aufgabe}

%% TODO Titel: Ich-kann-Satz fehlt in der Bank (Platzhalter: Kettenname) – Nr. 18
%% TODO Anweisung: ein Satz über der ersten Teilaufgabe fehlt in der Bank – Nr. 18
\begin{aufgabe}{Fläche oder Umfang: die passende Formel wählen}
\begin{teile}
\teil Um ein rundes Beet mit dem Radius $3$ m kommt eine Einfassung aus Metallband. Wie viel Band braucht man? Runde auf eine Stelle. \leerfeld
\end{teile}
\end{aufgabe}

%% TODO Titel: Ich-kann-Satz fehlt in der Bank (Platzhalter: Kettenname) – Nr. 19
%% TODO Anweisung: ein Satz über der ersten Teilaufgabe fehlt in der Bank – Nr. 19
\begin{aufgabe}{Fläche – Fehler finden, Begründen, Anwendung, Darstellungswechsel}
\begin{teile}
\teil Ein Kreis hat den Durchmesser $d = 12$ cm. Max rechnet die Fläche: \rechnung{A &= \pi \cdot 12^2 \\ A &\approx 452{,}4 \text{ cm}^2} Wo ist der Fehler?
\teil Ein Kreis wird in viele gleich große Tortenstücke zerschnitten. Die Stücke werden abwechselnd mit der Spitze nach oben und nach unten aneinandergelegt; es entsteht fast ein Rechteck. Begründe, warum die lange Seite ungefähr halb so lang ist wie der Umfang des Kreises.
\teil Eine Pizzeria bietet eine große Pizza ($36$ cm Durchmesser) für $12$ € an oder zwei kleine ($26$ cm Durchmesser) für zusammen $15$ €. Wo bekommt man mehr Pizza für einen Euro?
\teil Zu welcher Figur gehört der Flächenterm $\frac{1}{2} \cdot \pi \cdot 3^2$? Beschreibe sie mit Namen und Radius.
\end{teile}
\end{aufgabe}
